% Auto-gerado por C16_distorcao_conditional.py
\begin{table}[htbp]
\centering
\caption{Taxas de fluxo condicionais \`a defasagem idade-s\'erie no in\'icio do painel, Brasil, m\'edia 2018--2023 (ex-COVID).}
\label{tab:distorcao}
\begin{threeparttable}
\small
\begin{tabular}{lrrrrrr}
\toprule
& \multicolumn{4}{c}{Taxas (\%)} & & \\
\cmidrule(lr){2-5}
Subgrupo & Promo\c{c}\~ao & Repet\^encia & Evas\~ao & Migra\c{c}\~ao EJA & Soma & N \\
\midrule
\multicolumn{7}{l}{\textit{EF iniciais}} \\
\midrule
Em dia & 79.5 & 18.7 & 0.1 & 0.0 & 98.3 & 38.619 \\
Defasagem $\geq 2$ anos & 73.0 & 20.3 & 3.3 & 1.3 & 97.9 & 3.653 \\
\\[0.3em]
\multicolumn{7}{l}{\textit{EF finais}} \\
\midrule
Em dia & 77.6 & 20.1 & 0.4 & 0.2 & 98.3 & 29.346 \\
Defasagem $\geq 2$ anos & 63.2 & 22.0 & 7.3 & 4.6 & 97.1 & 6.028 \\
\\[0.3em]
\multicolumn{7}{l}{\textit{EM}} \\
\midrule
Em dia & 51.2 & 19.5 & 19.6 & 0.3 & 90.6 & 19.097 \\
Defasagem $\geq 2$ anos & 34.0 & 18.6 & 37.1 & 4.9 & 94.6 & 4.282 \\
\bottomrule
\end{tabular}
\begin{tablenotes}\footnotesize
\item \textit{Fonte:} Painel longitudinal da PNADC.
\item \textit{Notas:} Defasagem idade-s\'erie computada como diferen\c{c}a entre idade observada e idade-padr\~ao (entrada aos 6 anos no 1\textsuperscript{o} EF). Defasagem $\geq 2$ anos: idade $\geq$ idade-padr\~ao $+\,2$.
\end{tablenotes}
\end{threeparttable}
\end{table}
